\documentclass[../../main.tex]{subfiles}

\begin{document}
\sloppy

\chapter{Modern methods for correlated electron systems and superconductivity}\label{ch:modern_methods}

\section{Dynamical mean-field theory}\label{sec:dmft}

Dynamical mean-field theory is based on the idea that in the limit of infinite spatial dimension ($d\to\infty$), the self-energy $\Sigma_{\mathfrak{12}}^{\k}$ becomes \textit{purely local} \cite{Metzner1989}, where the Kronecker delta reduces the non-local self-energy to a purely local one\footnote{Due to the spatial locality of the self-energy, a Fourier transform to momentum space results in a flat dispersion of $\Sigma_{\text{DMFT}}^{\vv{k}}$.}. In high dimensions, the hopping amplitude must scale appropriately to remain finite. As a result, non-local correlations between lattice sites $i$ and $j$ are suppressed\footnote{Indeed, one can show that the diagrams contributing to the non-local character of the self-energy are suppressed by $\frac{1}{\sqrt{Z_{ij}}}$, where $Z_{ij}$ is the number of neighbors for a given distance $||\vv{R}_i-\vv{R}_j||$ from lattice site $i$.}. Starting from this discovery, it was shown \cite{Georges1992} that there is a connection between an auxiliary single-impurity Anderson model (SIAM) and the Hubbard model in infinite dimensions, which allowed for a derivation of a self-consistency cycle determining the purely local self-energy. This self-consistency can be calculated using approximate\footnote{e.g.\ self-consistent perturbation theory, iterated perturbation theory or non-crossing approximation} or numerically exact\footnote{e.g.\ quantum Monte Carlo (QMC), numerical renormalization group or even exact diagonalization} methods \cite{Held2007, Vollhardt2020}. A modern implementation of the DMFT procedure --- which is used in this thesis to calculate input data for the dynamical vertex approximation, more on that later --- is \texttt{w2dynamics} \cite{Wallerberger2019}, a continuous-time hybridization expansion (CT-HYB) QMC solver for the SIAM.
\\\\
Let us briefly outline the idea and self-consistency equations of DMFT in the following, see \cite{Held2007} for a more in-depth discussion. In high spatial dimensions, the number of neighboring sites $Z_{ij}$ grows as $Z_{ij}\sim d$. To ensure a well-defined competition between kinetic energy and Coulomb interaction, the hopping amplitudes must therefore scale as
\begin{align}
	t_{im,jn}\to t_{im,jn}^{*}=\frac{t_{im,jn}}{\sqrt{Z_{ij}}},
\end{align}
where $m,n$ denote source and target orbitals, respectively. Renormalizing the hopping elements automatically renormalizes the corresponding Green's function, which is written as
\begin{align}
	G_{\mathfrak{12}}^{\k}=\left [(i\nu+\mu)\delta_{\mathfrak{12}}-\varepsilon^{\vv{k}}_{\mathfrak{12}}-\Sigma_{\text{DMFT};\mathfrak{12}}^{\nu}\right ]^{-1},
\end{align}
where the Fourier transforms of the hopping amplitudes $\varepsilon^{\vv{k}}_{\mathfrak{12}}$ enter the Green's function, hence it scales similarly.

Diagrammatically, DMFT includes the local contribution of all topologically inequivalent Feynman diagrams which can also be obtained via the SIAM, however only if the interaction of the SIAM Hamiltonian has the same form as the original Hubbard Hamiltonian, 
\begin{align}
	\hat{\mathcal{H}}_{\text{AIM}}=\sum_{\mathfrak{ab};\vv{k}}\varepsilon^{\vv{k}}_{\mathfrak{ab}}\hat c^{\dagger}_{\vv{k};\mathfrak{a}}\hat c_{\vv{k};\mathfrak{b}}+\sum_{\mathfrak{ab}}\varepsilon^d_{\mathfrak{ab}}\hat d^\dagger_{\mathfrak{a}}\hat d_{\mathfrak{b}}+\frac12\sum_{\mathfrak{abcd}}U_{\mathfrak{abcd}}\hat d_{\mathfrak{a}}^{\dagger}\hat d_{\mathfrak{c}}^{\dagger} \hat d_{\mathfrak{d}} \hat d_{\mathfrak{b}}+\sum_{\mathfrak{ab};\vv{k}} \left [V^{\vv{k}}_{\mathfrak{ab}} \hat d_{\mathfrak{a}}^{\dagger}\hat c_{\vv{k};\mathfrak{b}}+\text{h.c.}\right ].
\end{align}
Here, $\hat d^{\dagger}_{\mathfrak{a}}$ ($\hat d_{\mathfrak{a}}$) creates (annihilates) an electron on the impurity, whereas the operators $\hat c^{\dagger}_{\vv{k};\mathfrak{a}}$ ($\hat c_{\vv{k};\mathfrak{a}}$) create (annihilate) bath --- i.e.\ non-interacting conduction --- electrons, which hybridize with the localized impurity electrons via $V^{\vv{k}}_{\mathfrak{ab}}$. For a successful mapping of the Hubbard model onto the AIM, there are three distinct equations in DMFT that describe a self-consistency cycle \cite{Held2007}. For a successful mapping of the Hubbard model, the impurity Green's function must equal the local Green's function of the corresponding lattice model. Since the local Green's function is not known a priori, the problem is solved self-consistently. One can start by \say{guessing} a local self-energy as a starting point and then iterate through the following set of equations until convergence:
\begin{subequations}
\begin{align}
	G^{\nu}_{\mathfrak{12}}&=\frac{1}{V_{\text{BZ}}}\integral{\text{BZ}}{}{\vv{k}}\left [(i\nu+\mu)\delta_{\mathfrak{12}}-\varepsilon^{\vv{k}}_{\mathfrak{12}}-\Sigma_{\text{DMFT};\mathfrak{12}}^{\nu}\right ]^{-1},\\
	\mathcal{G}^{\nu}_{0;\mathfrak{12}}&=\left [\big (G^{\nu}_{\mathfrak{12}}\big )^{-1}+\Sigma_{\text{DMFT};\mathfrak{12}}^{\nu}\right ]^{-1},\\
	G^{\nu}_{\mathfrak{12}}&=\mean{\mathcal{T}\hat d_{\mathfrak{1}}\hat d^{\dagger}_{\mathfrak{2}}}\qqand\label{eq:aim_gf}\\
	\Sigma_{\text{DMFT};\mathfrak{12}}^{\nu}&=\left (\mathcal{G}^{\nu}_{0;\mathfrak{12}}\right )^{-1}-\big (G^{\nu}_{\mathfrak{12}}\big )^{-1}.
\end{align}
\end{subequations}
Here, $\mathcal{G}^{\nu}_{0;\mathfrak{12}}$ denotes the non-interacting Green's function of the impurity. The numerically challenging part is the computation of the AIM's Green's function in \eqref{eq:aim_gf} for which a number of different solvers are available, see the beginning of \secref{sec:dmft}.

\section{Dynamical vertex approximation}\label{sec:dga}

DMFT --- as a local framework --- is very capable of describing strong quasiparticle renormalization and the Mott-Hubbard metal-to-insulator transition in heavy fermion systems and transition metal oxides accurately. Its capabilities extend way beyond that, where successful descriptions of (anti-)ferromagnetic structure effects including metamagnetism, influence on ferromagnetic properties through the lattice and Hund's coupling, scattering at spins or impurities and many more have been reported \cite{Held2007, Kotliar2006}. As broad as the range of DMFT's capabilities is, it comes with a major limitation: by construction, it neglects non-local correlations. Non-local correlations are responsible for numerous physical phenomena, such as the formation of valence bonds, pseudogaps, and d-wave superconductivity, as well as (para-)magnons and quantum critical behavior near phase transitions. These effects are particularly important in quasi-two-dimensional materials, where the infinite-dimensional limit of DMFT is far from realistic and a proper treatment of momentum or spatial structure is necessary to capture, for example, d-wave pairing.
\\\\
There are two main routes to extend DMFT beyond the local scope to also include non-local correlations \cite{Rohringer2012} (i) cluster extensions and (ii) diagrammatic extensions. Cluster extensions make use of a small cluster of impurities that are embedded in a bath of non-interacting conduction electrons \cite{Slezak2009}. Non-local interaction effects within this cluster are then taken into account, making it possible to capture short-range fluctuations. Diagrammatic extensions on the other hand extend DMFT by explicitly adding non-local diagrammatic contributions. The dynamical vertex approximation falls into the latter one and will be in the main focus of this thesis, more specifically the self-consistent \say{ladder} variant \cite{Kaufmann2021}. For other flavors of $\dga$, we refer the reader to Refs. \cite{DelRe2021, Valli2015, Schafer2021, Katanin2009, Held2008, Stobbe2022, Augustinsky2011, Tam2013}. For a first introduction to the \say{standard} description of the dynamical vertex approximation, we refer the reader to Ref.~\cite{Toschi2007}. 
\\\\
Here, we will extend the original description by including the purely non-local Coulomb interaction to the irreducible vertex and follow the derivation outlined in Ref.~\cite{Galler2017} closely. For completeness, we will also write down the set of equations analogously to \cite{Worm2023a} for the case of non-local interactions. For clarity, we will govern those two separate derivation branches each their own section. The reason we include both derivations in this thesis is simple: Previously, we have written down two distinct ways to treat explicit vertex asymptotics. The first one, where explicit kernel-functions are used, see \secref{sec:vertex_asymptotics_kernel}, works well if one explicitly employs the full vertex $F$ in the equations, as \cite{Galler2017} does. However, the second derivation, where we do not need the full vertex in the equations explicitly, benefits when treating vertex asymptotics with the \say{$U$-range} approach, see \secref{sec:vertex_asymptotics_urange}, since the numerical implementation of these asymptotics are more resourceful and one does not need local quantities on a very large frequency grid, i.e.\ it is easier to increase the size of the asymptotic range compared to the kernel-function approach. We realized that the second approach is easier to realize numerically and the code developed in this thesis employs the \say{$U$-range} asymptotics as formulated by \cite{Kitatani2019}, which is also tested against the implementation of \cite{Worm2023}.
\\\\
The $\dga$ uses a converged DMFT solution, thus including all local diagrams. It then calculates a non-local approximation of the vertex, which in turn is inserted into the Schwinger-Dyson equation (\ref{eq:sde}) to obtain a new self-energy. We will present the \say{ladder}-approach, which reduces the numerical complexity by assuming the (ph, $\phbar$ and fully) irreducible vertices to be purely local and equal to those of DMFT. Furthermore, it also assumes pure locality of pp-reducible diagrams\footnote{This is not a priori justified, however, results of the full (parquet) $\dga$ equations yield a negligible contribution of the particle-particle channel to the full vertex $F$ in the presence of strong spin fluctuations close to the antiferromagnetic phase transition in two dimensions \cite{Held2008}. It is therefore reasonable to neglect the non-locality of the particle-particle channel diagrams. For superconducting properties on the other hand, these non-local diagrams are often essential.}. For the following, we assume that we already have an existing converged DMFT solution at hand, containing all local diagrams. The quantities we require for the execution of the dynamical vertex approximation are the DMFT self-energy and the local two-particle Green's function in the particle-hole channel\footnote{If one would use the asymptotics of the vertex functions as described in \secref{sec:vertex_asymptotics_kernel}, one also requires the local two-particle Green's function in the $\pp$-channel.}.

At the heart of $\dga$ is the irreducible vertex in the particle-hole channel\footnote{From now on, if not explicitly specified otherwise, we will always assume the ph-frequency notation and channel reducibility of $\Gamma$ and $\Phi$ to be particle-hole.}, which we extend by the purely non-local Coulomb interaction,
\begin{align}\label{eq:non_local_gamma_abinitio_dga}
	\Gamma^{\q\k\kp}_{\mathfrak{1234};\ss'}=\Gamma^{\omega\nu\nu'}_{\mathfrak{1234};\ss'}+\mathcal{V}^{\vv{qkk}'}_{\mathfrak{1234};\ss'},
\end{align}
where $\mathcal{V}^{\vv{qkk}'}_{\mathfrak{1234};\ss'}$ is given by the crossing-symmetric expression (similar to \eqref{eq:crossing_symmetric_uq})
\begin{align}
	\mathcal{V}^{\vv{qkk}'}_{\mathfrak{1234};\ss'}=V^{\vv{q}}_{\mathfrak{1234}}-\delta_{\ss'}V^{\vv{k}'-\vv{k}}_{\mathfrak{1432}}.
\end{align}
The irreducible vertex can now be seen as containing all local irreducible diagrams plus the non-locality of the first order, providing a natural extension to this framework. This means that we have now included the first-order term --- which is also the major contributor to the irreducible vertex --- fully momentum-dependent. We will, similar to Ref.~\cite{Galler2017}, exclude from now on the explicit transversal particle-hole contribution to the non-local irreducible vertex in \eqref{eq:non_local_gamma_abinitio_dga}, similar to the GW approach \cite{Tomczak2017}, which will reduce the complexity of the $\dga$ equations considerably, since two out of three momentum indices vanish; the only leftover one from now on is the transfer momentum $\vv{q}$. This yields for the irreducible vertex in $\ph$-channel for the $\dens/\magn$ spin combinations
\begin{subequations}\label{eq:gamma_dens_magn_dga}
\begin{align}
	\Gamma^{\q\nu\nu'}_{\dens;\mathfrak{1234}}&=\Gamma^{\omega\nu\nu'}_{\dens;\mathfrak{1234}}+2 V^{\vv{q}}_{\mathfrak{1234}}\qqand\\
	\Gamma^{\q\nu\nu'}_{\magn;\mathfrak{1234}}&=\Gamma^{\omega\nu\nu'}_{\magn;\mathfrak{1234}}.\label{eq:gamma_magn_dga}
\end{align}
\end{subequations}
From this, the Bethe-Salpeter equation in the $\ph$-channel of \eqref{eq:bethe_salpeter_vertex_three_channels} reduces to
\begin{align}
	F_{\dens/\magn;\mathfrak{1234}}^{\q\nu\nu'}=\Gamma^{\q\nu\nu'}_{\dens/\magn;\mathfrak{1234}}-\frac1\beta\sum_{\mathfrak{abcd};\nu_1}\Gamma^{\q\nu\nu_1}_{\dens/\magn;\mathfrak{12ab}}\chi_{0;\mathfrak{badc}}^{\q\nu_1\nu_1} F_{\dens/\magn;\mathfrak{cd34}}^{\q\nu_1\nu'},
\end{align}
where
\begin{align}
	\chi_{0;\mathfrak{1234}}^{\q\nu\nu'}=\sum_{\vv{kk}'}\chi_{0;\mathfrak{1234}}^{\q\k\kp}=-\beta\sum_{\vv{kk}'}G^{\k}_{\mathfrak{14}}G^{\kp-\q}_{\mathfrak{32}}.
\end{align}
From now on we will sometimes work with grouped indices, as described at the end of \secref{sec:bethe_salpeter}, to simplify the following derivation of the $\dga$ equations. Notice that the grouped indices are now $\{ n_1,n_2,\nu \}$ and $\{ n_3,n_4,\nu' \}$, respectively, and do not include any momenta anymore. We will denote quantities written in grouped indices as bold quantities, similar to before. We can now see that the numerical implementation of the $\dga$ equations is easier, since the vertex matrices decrease in size from $(N_{\vv{k}}\cdot N_{\nu} \cdot N_{\text{bands}}^2)\times(N_{\vv{k}}\cdot N_{\nu} \cdot N_{\text{bands}}^2) \to (N_{\nu} \cdot N_{\text{bands}}^2)\times(N_{\nu} \cdot N_{\text{bands}}^2)$, which reduces the storage capacity by $N_{\vv{k}}^2$ and the number of operations by $N_{\vv{k}}^2$ for matrix multiplication and $N_{\vv{k}}^3$ for matrix inversion compared to the full parquet approach.

In principle, it would be straightforward to retrieve the full vertex via $\bm F_{\r}^{\q}=\left [(\bm\Gamma_{\r}^{\q})^{-1}+\bm\chi_{0}^{\q}\right ]^{-1}$ from \eqref{eq:bse_matrix}. This poses numerical problems, as already stated at the end of \secref{sec:schwinger_dyson}. Thus, we rewrite the inversion of the irreducible vertex $\bm\Gamma$ using the local variant of \eqref{eq:bethe_salpeter_vertex_three_channels} and \eqref{eq:gamma_dens_magn_dga} and after some algebra, the full vertex in the density and magnetic spin combinations reads \cite{Galler2017}
\begin{subequations}
\begin{align}
\begin{split}
	\bm F_{\dens}^{\q}&=\big [\bm F_{\dens}^{\omega}+2\bm V^{\vv{q}}\left (1-\bm\chi_{0}^{\omega}\bm F_{\dens}^{\omega}\right )\big ]\!\left [\bm\unity+\bm\chi_{0}^{\nl;\q}\bm F_{\dens}^{\omega}-2\bm\chi_{0}^{\q}\bm V^{\vv{q}}\left (1-\bm\chi_{0}^{\omega}\bm F_{\dens}^{\omega}\right )\right ]^{-1}
\end{split}\;\quad\text{and}\quad\\
	\bm F_{\magn}^{\q}&=\bm F_{\magn}^{\omega}\left [\bm\unity+\bm\chi_{0}^{\nl;\q}\bm F_{\magn}^{\omega}\right ]^{-1}.
\end{align}
\end{subequations}
Here, the purely non-local susceptibility $\bm\chi_{0}^{\nl;\q}$ is given by
\begin{align}
	\bm\chi_{0}^{\nl;\q}=\bm\chi_{0}^{\q}-\bm\chi_{0}^{\omega}.
\end{align}
Rewriting the full vertex this way avoids the aforementioned numerical difficulties that lie in the inversion of the irreducible vertex $\bm\Gamma$, however we still have to deal with the inclusion of the transversal particle-hole channel, since the above equations were generated though the BSE (\ref{eq:bethe_salpeter_vertex_three_channels}) in the $\ph$-channel only. The BSE only carries over the swapping symmetry of the two-particle Green's function to the (ir-)reducible vertices $\bm\Gamma$ and $\bm\Phi$. The crossing symmetry of the full vertex can be restored by explicitly adding diagrams that belong to the transversal particle-hole channel. We therefore modify the parquet-decomposition in \eqref{eq:parquet_decomposition} in the following way
\begin{align}\label{eq:full_vertex_crossing_symmetric_dga}
	F^{\q\k\kp}_{\dens;\mathfrak{1234}}&=F^{\omega\nu\nu'}_{\dens;\mathfrak{1234}}+2 V_{\mathfrak{1234}}^{\vv{q}}+\left (\Phi^{\q\nu\nu'}_{\dens;\mathfrak{1234}}-\Phi^{\omega\nu\nu'}_{\dens;\mathfrak{1234}}\right )+\left (\Phi^{\q\k\kp}_{\dens;\phbar;\mathfrak{1234}}-\Phi^{\omega\nu\nu'}_{\dens;\phbar;\mathfrak{1234}}\right )\qqand\\
	F^{\q\k\kp}_{\magn;\mathfrak{1234}}&=F^{\omega\nu\nu'}_{\r;\mathfrak{1234}}+\left (\Phi^{\q\nu\nu'}_{\magn;\mathfrak{1234}}-\Phi^{\omega\nu\nu'}_{\magn;\mathfrak{1234}}\right )+\left (\Phi^{\q\k\kp}_{\magn;\phbar;\mathfrak{1234}}-\Phi^{\omega\nu\nu'}_{\magn;\phbar;\mathfrak{1234}}\right ).
\end{align}
Here, all diagrams in the $\pp$-channel and all fully irreducible diagrams (except the bare interaction $V^{\vv{q}}_{r;\mathfrak{1234}}$) are local. We furthermore find via the local and non-local BSE,
\begin{subequations}
\begin{align}
	\left (\Phi^{\q\nu\nu'}_{\dens;\mathfrak{1234}}-\Phi^{\omega\nu\nu'}_{\dens;\mathfrak{1234}}\right )&=F_{\dens;\mathfrak{1234}}^{\nl;\q\nu\nu'}-2 V^{\vv{q}}_{\mathfrak{1234}}\qqand\\
	\left (\Phi^{\q\nu\nu'}_{\magn;\mathfrak{1234}}-\Phi^{\omega\nu\nu'}_{\magn;\mathfrak{1234}}\right )&=F_{\magn;\mathfrak{1234}}^{\nl;\q\nu\nu'},
\end{align}
\end{subequations}
where $F_{\r;\mathfrak{1234}}^{\nl;\q\nu\nu'}$ is defined analogously to the purely non-local susceptibility 
\begin{align}
	F_{\r;\mathfrak{1234}}^{\nl;\q\nu\nu'}=F_{\r;\mathfrak{1234}}^{\q\nu\nu'}-F_{\r;\mathfrak{1234}}^{\omega\nu\nu'}.
\end{align}
We can express the difference of the transversal particle-hole diagrams of \eqref{eq:full_vertex_crossing_symmetric_dga} with the help of  \eqref{eq:phi_ph_bar_dens_magn}, yielding
\begin{subequations}
\begin{align}
	\left (\Phi^{\q\k\kp}_{\dens;\phbar;\mathfrak{1234}}-\Phi^{\omega\nu\nu'}_{\dens;\phbar;\mathfrak{1234}}\right )&=-\frac12 F_{\dens;\mathfrak{3214}}^{\nl;(\kp-\k)(\nu'-\omega)\nu'}-\frac32 F_{\magn;\mathfrak{3214}}^{\nl;(\kp-\k)(\nu'-\omega)\nu'}\qqand\\
	\left (\Phi^{\q\k\kp}_{\magn;\phbar;\mathfrak{1234}}-\Phi^{\omega\nu\nu'}_{\magn;\phbar;\mathfrak{1234}}\right )&=-\frac12 F_{\magn;\mathfrak{3214}}^{\nl;(\kp-\k)(\nu'-\omega)\nu'}+\frac12 F_{\magn;\mathfrak{3214}}^{\nl;(\kp-\k)(\nu'-\omega)\nu'}.
\end{align}
\end{subequations}
We can combine the results of \eqref{eq:full_vertex_crossing_symmetric_dga} which now states
\begin{subequations}\label{eq:full_vertex_ladder_dga_crossing_symmetric}
\begin{align}
	F^{\q\k\kp}_{\dens;\mathfrak{1234}}&=F^{\omega\nu\nu'}_{\dens;\mathfrak{1234}}\,+F^{\nl;\q\nu\nu'}_{\dens;\mathfrak{1234}}-\frac12 F_{\dens;\mathfrak{3214}}^{\nl;(\kp-\k)(\nu'-\omega)\nu'}-\frac32 F_{\magn;\mathfrak{3214}}^{\nl;(\kp-\k)(\nu'-\omega)\nu'}\qqand\\
	F^{\q\k\kp}_{\magn;\mathfrak{1234}}&=F^{\omega\nu\nu'}_{\magn;\mathfrak{1234}}+F^{\nl;\q\nu\nu'}_{\magn;\mathfrak{1234}}-\frac12 F_{\dens;\mathfrak{3214}}^{\nl;(\kp-\k)(\nu'-\omega)\nu'}+\frac12 F_{\magn;\mathfrak{3214}}^{\nl;(\kp-\k)(\nu'-\omega)\nu'}.
\end{align}
\end{subequations}
Within the approximations introduced above, the connected part of the self-energy (see \eqref{eq:sde} and below) loses the momentum-dependent contribution from the non-local bare interaction $V^{\vv{k}'-\vv{k}}_{\mathfrak{1234}}$ and now reads
\begin{align}\label{eq:sde_with_phbar}
	\Sigma^{\k}_{\mathfrak{12}}&=\Sigma_{\text{HF};\mathfrak{12}}^{\vv{k}} -\frac1\beta\sum_{\substack{\mathfrak{abcdef}\\\q\nu'}}\left [\mathcal{U}^{\vv{q}}_{\mathfrak{a1bc}}\chi^{\q\nu'\nu'}_{0;\mathfrak{cbed}} F^{\q\nu'\nu}_{\dens;\mathfrak{de2f}}-\frac12\tilde{\mathcal{U}}_{\mathfrak{a1bc}}\chi^{\q\nu'\nu'}_{0;\mathfrak{cbed}}\left [F^{\nl;\q\nu'\nu}_{\dens;\mathfrak{2edf}}+3 F^{\nl;\q\nu'\nu}_{\magn;\mathfrak{2edf}}\right ]\right ]G_{\mathfrak{af}}^{\k-\q}.
\end{align}
It is now time to split up the derivation in two paths, one devoted to the steps outlined in \cite{Galler2017}, coined \say{Approach 1}, and one devoted to the steps outlined in \cite{Worm2023a}, coined \say{Approach 2}.
\\\\
\textbf{Approach 1} --- From the purely local part of \eqref{eq:sde_with_phbar}, which is contained in the first term in the brackets, one can extract the DMFT self-energy, since the $\ph$-channel contribution with the local $U_{\mathfrak{1234}}$ can be rewritten as
\begin{align}
	\Sigma_{\ph;U;\mathfrak{12}}^{\k}&=\Sigma_{\text{HF};\mathfrak{12}}^{\vv{k}}-\frac{1}{\beta}\sum_{\q\nu'}U_{\mathfrak{a1bc}}\left (\chi^{\nl;\q\nu'\nu'}_{0;\mathfrak{cbed}}+\chi^{\omega\nu'\nu'}_{0;\mathfrak{cbed}}\right ) F^{\omega\nu'\nu}_{\dens;\mathfrak{de2f}} G_{\mathfrak{af}}^{\k-\q}\\
	&=\Sigma_{\text{DMFT};\mathfrak{12}}^{\nu}+\Sigma_{\text{HF};\mathfrak{12}}^{\nl;\vv{k}}-\frac{1}{\beta}\sum_{\q\nu'}U_{\mathfrak{a1bc}}\chi^{\nl;\q\nu'\nu'}_{0;\mathfrak{cbed}}F^{\omega\nu'\nu}_{\dens;\mathfrak{de2f}} G_{\mathfrak{af}}^{\k-\q},
\end{align}
where $\Sigma_{\text{HF};\mathfrak{12}}^{\nl;\vv{k}}$ is the purely non-local Hartree-Fock term with the contributions from $V^{\vv{q}}_{\mathfrak{1234}}$ but not $U_{\mathfrak{1234}}$, see \eqref{eq:sde}. We can take advantage of the momentum and frequency structure in the SDE above to simplify the numerical calculation of the self-energy. Therefore, we --- in analogy to Ref.~\cite{Galler2017} --- define three Hedin vertices in $\r\in\{\dens,\magn\}$ spin combinations,
\begin{align}\label{eq:three_leg_vertices_abinitio_dga}
	\zeta_{\r;\mathfrak{1234}}^{\omega\nu}&=\sum_{\mathfrak{ab;\nu'}}\chi_{0;\mathfrak{12ab}}^{\omega\nu'\nu'}F_{\r;\mathfrak{ba34}}^{\omega\nu'\nu},\\
	\zeta_{\r;\mathfrak{1234}}^{\q\nu}&=\sum_{\mathfrak{ab;\nu'}}\chi_{0;\mathfrak{12ab}}^{\nl;\q\nu'\nu'}F_{\r;\mathfrak{ba34}}^{\omega\nu'\nu}\qqand\\
	\xi_{\r;\mathfrak{1234}}^{\q\nu}&=\sum_{\mathfrak{ab;\nu'}}\left [\chi_{0;\mathfrak{12ab}}^{\q\nu'\nu'}F_{\r;\mathfrak{ba34}}^{\q\nu'\nu}-\chi_{0;\mathfrak{12ab}}^{\omega\nu'\nu'}F_{\r;\mathfrak{ba34}}^{\omega\nu'\nu}\right ].
\end{align}
The completely local Hedin vertex can straightforwardly be retrieved from DMFT. The vertex $\xi_{\mathfrak{1234}}^{\q\nu}$ in spin combinations $\dens/\magn$ can be calculated efficiently with a matrix inversion
\begin{subequations}
\begin{align}
	\bm\xi_{\dens}^{\q}&=\left (\vec{\bm\unity}+\bm\zeta_{\dens}^{\omega}\right )\left (\left [\bm\unity-\bm\chi_{0}^{\nl;\q}\bm F_{\dens}^{\omega}-2\bm\chi_{0}^{\q}\bm V^{\vv{q}}\left (\vec{\bm\unity}+\bm\zeta_{\dens}^{\omega}\right )\right ]^{-1}-\bm\unity\right )\qqand\\
	\bm\xi_{\magn}^{\q}&=\left (\vec{\bm\unity}+\bm\zeta_{\magn}^{\omega}\right )\left (\left [\bm\unity-\bm\chi_{0}^{\nl;\q}\bm F_{\magn}^{\omega}\right ]^{-1}-\bm\unity\right ),
\end{align}
\end{subequations}
where $\vec{\bm\unity}_{\mathfrak{1234}}=\delta_{\mathfrak{14}}\delta_{\mathfrak{23}}$, contrary to $\bm\unity$, which is given by $\bm\unity_{\mathfrak{1234}}^{\nu\nu'}=\delta_{\mathfrak{14}}\delta_{\mathfrak{23}}\delta_{\nu\nu'}$. Combining these expressions and simplifying it \cite{Galler2017} yields the $\dga$ self-energy
\begin{align}\label{eq:sde_abinitio_dga}
\begin{split}
	\bm\Sigma_{\text{D}\Gamma\text{A}}^{\k}&=\bm\Sigma_{\text{DMFT}}^{\nu}+\bm\Sigma_{\text{HF}}^{\nl;\vv{k}}-\frac1\beta\sum_{\q}\left (\bm U + \bm V^{\vv{q}}-\frac{\tilde{\bm U}}{2}\right )\bm\xi_{\dens}^{\q}\bm G^{\k-\q}\\
	&\qquad+\frac{3}{2\beta}\sum_{\q}\tilde{\bm U}\bm\xi_{\magn}^{q}\bm G^{\k-\q}-\frac1\beta\sum_{\q}\big (\bm V^{\vv{q}}\bm\zeta_{\dens}^{\omega}-\bm U\bm\zeta_{\dens}^{\q}\big )\bm G^{\k-\q}.
\end{split}
\end{align}
\\\\
\textbf{Approach 2} --- For an alternative formulation using the three-leg vertices already derived in \secref{sec:schwinger_dyson}, see \eqref{eq:vrg}, we have to rewrite \eqref{eq:sde_with_phbar} with the full ladder vertices $F^{\q\nu\nu'}_{\dens/\magn;\mathfrak{1234}}$,
\begin{align}
	F^{\nl;\q\nu'\nu}_{\dens/\magn;\mathfrak{1234}}=F^{\q\nu'\nu}_{\dens/\magn;\mathfrak{1234}}-F^{\omega\nu'\nu}_{\dens/\magn;\mathfrak{1234}},
\end{align}
such that
\begin{align}
	\Sigma^{\k}_{\mathfrak{12}}&=\Sigma_{\text{HF};\mathfrak{12}}^{\vv{k}} -\frac1\beta\sum_{\substack{\mathfrak{abcdef}\\\q\nu'}}\left [\mathcal{U}^{\vv{q}}_{\mathfrak{a1bc}}\chi^{\q\nu'\nu'}_{0;\mathfrak{cbed}} F^{\q\nu'\nu}_{\dens;\mathfrak{de2f}}-\frac12\tilde{\mathcal{U}}_{\mathfrak{a1bc}}\chi^{\q\nu'\nu'}_{0;\mathfrak{cbed}}\left [F^{\q\nu'\nu}_{\dens;\mathfrak{2edf}}+3 F^{\q\nu'\nu}_{\magn;\mathfrak{2edf}}\right ]\right ]G_{\mathfrak{af}}^{\k-\q} - \Sigma^{\k}_{\text{dc};\mathfrak{12}},
\end{align}
where $\Sigma^{\k}_{\text{dc};\mathfrak{12}}$, which is merely a double-counting correction, is given by
\begin{align}\label{eq:sigma_double_counting}
	\Sigma^{\k}_{\text{dc};\mathfrak{12}} = \frac{1}{2\beta}\sum_{\substack{\mathfrak{abcdef}\\\q\nu'}}\tilde{\mathcal{U}}_{\mathfrak{a1bc}}\chi^{\q\nu'\nu'}_{0;\mathfrak{cbed}}\left [F^{\omega\nu'\nu}_{\dens;\mathfrak{2edf}}+3 F^{\omega\nu'\nu}_{\magn;\mathfrak{2edf}}\right ]G_{\mathfrak{af}}^{\k-\q}.
\end{align}
With the definition of the three-leg vertex from \eqref{eq:vrg}, we can write the upper equation as
\begin{align}\label{eq:sde_vrg_ladder_full}
	\Sigma^{\k}_{\dga;\mathfrak{12}}&=\Sigma_{\text{HF};\mathfrak{12}}^{\vv{k}}-\frac12\sum_{\substack{\q;\mathfrak{abcd}\\\mathfrak{efgh}}}\Bigg [\,\mathcal{U}^{\vv{q}}_{\dens;\mathfrak{a1bc}}\gamma_{\dens;\mathfrak{cbed}}^{\q\nu}\!\left (\unity_{\mathfrak{de2h}}-\mathcal{U}^{\vv{q}}_{\dens;\mathfrak{defg}}\chi^{\q}_{\dens;\mathfrak{gf2h}}\right )\nonumber\\
	&\qquad\qquad\quad+3\, \mathcal{U}^{\vv{q}}_{\magn;\mathfrak{a1bc}}\gamma_{\magn;\mathfrak{cbed}}^{\q\nu}\!\left (\unity_{\mathfrak{de2h}}-\mathcal{U}^{\vv{q}}_{\magn;\mathfrak{defg}}\chi^{\q}_{\magn;\mathfrak{gf2h}}\right )-2\,\mathcal{U}^{\vv{q}}_{\magn;\mathfrak{a12h}}\Bigg]G_{\mathfrak{ah}}^{\k-\q} - \Sigma^{\k}_{\text{dc};\mathfrak{12}}\nonumber\\
\begin{split}
	&=\Sigma_{\text{HF};\mathfrak{12}}^{\vv{k}}-\frac12\sum_{\substack{\q;\mathfrak{abcd}\\\mathfrak{efgh}}}\Bigg [\bigg (2U_{\mathfrak{a1bc}} + 2V^{\vv{q}}_{\mathfrak{a1bc}} - U_{\mathfrak{acb1}} \bigg )\,\gamma_{\dens;\mathfrak{cbed}}^{\q\nu}\times\\
	&\qquad\qquad\qquad\qquad\times\Big (\unity_{\mathfrak{de2h}}-\left (2U_{\mathfrak{defg}}+2V^{\vv{q}}_{\mathfrak{defg}}-U_{\mathfrak{dgfe}}\right )\chi^{\q}_{\dens;\mathfrak{gf2h}}\Big )\\
	&\qquad\qquad\qquad\qquad-3 U_{\mathfrak{acb1}}\gamma_{\magn;\mathfrak{cbed}}^{\q\nu}\!\left (\unity_{\mathfrak{de2h}}+U_{\mathfrak{dgfe}}\chi^{\q}_{\magn;\mathfrak{gf2h}}\right )\!+2U_{\mathfrak{ah21}}\!\Bigg]G_{\mathfrak{ah}}^{\k-\q} - \Sigma^{\k}_{\text{dc};\mathfrak{12}}.
\end{split}
\end{align}

\section{Self-consistent dynamical vertex approximation}

The self-energy obtained from a simple \say{one-shot} calculation with the ladder-$\dga$ does not (always) show the correct asymptotic behavior. This deficiency is even more pronounced for higher values of the physical susceptibilities \cite{Kaufmann2021}. In addition, the susceptibilities that can be derived from the full vertex diverge at the DMFT N\'eel temperature. This is concerning, since it therefore violates the Mermin-Wagner theorem \cite{Mermin1966} for the case of two-dimensional systems. This can be resolved by applying a so-called $\lambda$-correction \cite{Moriya1985}, where enforcing a specific sum rule of the spin (and charge) susceptibilities yields a renormalization parameter $\lambda$, which is in turn used to renormalize physical susceptibilities in the $\dga$ routine.

The formalism for the $\lambda$-correction is currently only available in the one-band regime, an extension to multi-orbital systems has not been derived yet. The reasons being (i) $\lambda$ would become a matrix object with as many entries as spin-orbit combinations exist, yielding a multidimensional optimization problem; and (ii) the solution of this multidimensional problem can be non-unique and there are no rules of which of these solutions to choose. Hence the approach using a regularization scheme with the help of a $\lambda$-correction is not suitable for multi-orbital systems, which are in the focus of this thesis. Therefore, to \say{repair} the violation of the Mermin-Wagner theorem, we will employ a different scheme, a so-called self-consistency iteration \cite{Kaufmann2021}, more on that later. 

Since the $\lambda$-correction is currently not possible for multi-band systems, simply using a non-renormalized one-shot ladder-$\dga$ calculation might be a first approach. The issue is that it is not a viable description in parameter regimes where non-local correlations are large \cite{Kaufmann2021}, e.g.\ where the DMFT susceptibilities diverge due to a phase transition. The source of this is two-fold: (i) due to the assumed locality of the particle-particle reducible diagrams in the ladder-$\dga$ equations, which would dampen the particle-hole fluctuations \cite{Rohringer2016}, the ladder diagrams of the particle-hole channel lack their non-local contributions; and (ii) the self-energy entering the Bethe-Salpeter equations is only the local DMFT self-energy. Hence any non-local contributions due to the self-energy are simply missing. The latter problem can be solved by applying a self-consistency scheme within the $\dga$ equations as shown in \figref{fig:flowchart_self_consistency} \cite{Kaufmann2021}.
\begin{figure}[ht!]
	\centering
  	\includestandalone[scale=1]{Graphics/Diagrams/flowchart_self_consistency/flowchart_self_consistency}
  	\caption{This flowchart shows the $\dga$ self-consistency cycle. The first box is concerned with the bandstructure and interaction calculations which can be obtained from simple model studies or e.g.\ DFT or cRPA. The next box denotes the DMFT calculation which is performed using the previously acquired bandstructure and interaction. From the converged DMFT calculation we extract the DMFT self-energy and the local two-particle Green's function. The three boxes below describe the self-consistency cycle of the $\dga$ approach. The bandstructure, interaction and output quantities from DMFT, i.e.\ all purely local (ir-)reducible vertex functions are kept constant throughout the whole $\dga$ self-consistency cycle.}
  	\label{fig:flowchart_self_consistency}
\end{figure}
The steps to acquire a converged $\dga$ routine with an appropriate self-energy and vertex functions are as follows:
\begin{enumerate}
	\item Initially, the method requires a band structure $\varepsilon^{\vv{k}}_{\mathfrak{12}}$ and a (non-local) interaction $\mathcal{U}^{\vv{q}}_{\mathfrak{1234}}$, which can come from model studies or, for example, from DFT or cRPA. Using this momentum-dependent band structure and interaction, a DMFT calculation is performed with the chosen (imaginary-time) impurity solver. Once the DMFT calculation has converged, the necessary local one- and two-particle quantities --- such as the local self-energy, the local irreducible vertex, and the full vertex --- are extracted. Each of these local quantities remains fixed throughout the subsequent self-consistency cycle. This is particularly important for the irreducible vertex $\Gamma$ and the full vertex $F$, which are required for the Bethe-Salpeter equation and the double-counting correction to $\Sigma$, respectively. All of these quantities then serve as input for the following $\mathrm{D}\Gamma\mathrm{A}$ calculation.
	\item The first step in the sc-$\dga$ routine is the calculation of the one-particle Green's function via the Dyson equation (\ref{eq:dyson_eq}). This is done with the local self-energy from DMFT in the first iteration and then with the non-local self-energy calculated from the non-local SDE of the previous iteration.
	\item The next step is to update the chemical potential from the newly acquired Green's function to make sure the particle number stays constant. 
	\item We then proceed with the calculation of the necessary (non-local) three-leg vertices $\gamma^{\q\nu}_{\dens/\magn;\mathfrak{1234}}$, see \eqref{eq:vrg}, and the generalized physical susceptibilities $\chi^{\q}_{\dens/\magn;\mathfrak{1234}}$ from \eqref{eq:chi_phys_asympt}. While we keep $\Gamma^{\omega\nu\nu'}_{\dens/\magn;\mathfrak{1234}}$ constant throughout the whole process, $\gamma^{\q\nu}_{\dens/\magn;\mathfrak{1234}}$ and $\chi^{\q}_{\dens/\magn;\mathfrak{1234}}$ are changing each iteration due to the updated Green's function (and bare bubble) and its feedback into these non-local quantities.
	\item The self-energy of the current iteration is finally calculated via \eqref{eq:sde_abinitio_dga} using the previously acquired vertex functions and susceptibilities. To complete the cycle, the self-energy then enters the Green's function again via the Dyson equation, see step 2.
\end{enumerate}
To ensure a steady convergence, the self-energy that enters the Green's function is either mixed linearly with a parameter $\alpha$, such that $\Sigma_n=\alpha\Sigma_n + (1-\alpha)\Sigma_{n-1}$, where $n$ denotes the iteration number and $\alpha$ is typically chosen around $0.2-0.4$ or with other mixing schemes, such as Periodic Pulay mixing \cite{Banerjee2016}. This accelerated and predictive mixing scheme has been used previously in numerical calculations \cite{Peil2022} with significant improvement in convergence speed and behavior over the linear mixing strategy, despite the additional computational and storage cost. 

To determine whether the self-energy is converged, a simple criteria $|\Sigma_n^{\k}-\Sigma_{n-1}^{\k}|<\varepsilon$ is used. Here, $n$ denotes the current iteration and $\varepsilon$ is a small number of $\sim\mathcal{O}\left (10^{-4}\right )$.

From the one-particle Green's function, as well as the vertices and susceptibilities that are directly accessible at this stage, one can extract additional physical quantities, including the optical conductivity and superconducting properties. For instance, the critical temperature can be estimated and the symmetry of the superconducting gap can be determined by solving the Eliashberg equation, which is treated as a post-processing step described in the following section.

\section{Eliashberg equation}

After a successful converged ladder-$\dga$ routine, we have the most important non-local one- and two-particle quantities available. Starting from this, we can further study superconductivity with the so-called Eliashberg equations \cite{Eliashberg1960, Marsiglio2020}. In this context, we introduce, additionally to the one-particle Green's function of \eqref{eq:one_p_green_expectation_val}, two \textit{anomalous} Green's functions \cite{Kaser2022}
\begin{subequations}
\begin{align}
	F_{\mathfrak{12}}^{\k}&=-\mean{\mathcal{T}\left [\hat{c}_{\k;\mathfrak{1}}\hat{c}_{-\k;\mathfrak{2}}\right ]} \qqand\\
	\overline{F}_{\mathfrak{12}}^{\k}&=-\mean{\mathcal{T}\left [\hat{c}^{\dagger}_{\k;\mathfrak{1}}\hat{c}^{\dagger}_{-\k;\mathfrak{2}}\right ]}.
\end{align}
\end{subequations}
Here, $F_{\mathfrak{12}}^{\k}$ and $\overline{F}_{\mathfrak{12}}^{\k}$ are not to be confused with the full vertex $F^{\q\k\kp}_{\mathfrak{1234}}$. The anomalous Green's functions are rather unusual contractions of either two electrons or two holes, contrary to the regular Green's function, which is a contraction of an electron and a hole. These two contractions will vanish in the normal state because of particle conservation, however they are of finite value in the superconducting phase. 

At this point is is helpful to switch to the Nambu formalism \cite{Nambu1950}, where the basis is spanned by the spinors
\begin{align}
	\hat\psi_{\mathfrak{1}}^{\k}=\left (\hat{c}_{\k;\mathfrak{1};\uparrow}, \hat{c}_{\k;\mathfrak{1};\downarrow},\hat{c}^{\dagger}_{-\k;\mathfrak{1};\uparrow},\hat{c}^{\dagger}_{-\k;\mathfrak{1};\downarrow}\right )^{\text{T}}.
\end{align}
In this basis, the anomalous Green's functions are (block-)off-diagonal entries of the (generalized) Gro'kov single-particle propagator \cite{Yanase2003, Mineev1999}
\begin{align}
\begin{split}
	\hat{\bm G}_{\mathfrak{12}}^{\k}&=-\mean{\mathcal{T}\left [\hat{\psi}_{\mathfrak{1}}^{\k}\hat{\psi}_{\mathfrak{2}}^{\k;\dagger}\right ]}\\
	&=\begin{pmatrix}
		G_{\mathfrak{12}}^{\k} & -F_{\mathfrak{12}}^{\k}	\\
		-\overline{F}_{\mathfrak{12}}^{\k} & \overline{G}_{\mathfrak{12}}^{\k}
	\end{pmatrix},
\end{split}
\end{align}
where $G_{\mathfrak{12}}^{\k}$ and $F_{\mathfrak{12}}^{\k}$ are $2n \times 2n$ matrices in spin-orbital space, where $n$ denotes the number of bands. The diagonal entries of the Gro'kov single-particle propagator are the regular single-particle ($ G_{\mathfrak{12}}^{\k}$) and single-hole ($\overline{G}_{\mathfrak{12}}^{\k}=-G_{\mathfrak{12}}^{-\k}$) propagators. After a Fourier transform to Matsubara frequencies, the generalized Dyson equation can be written as follows:
\begin{align}\label{eq:grokov_dyson}
	\hat{\bm G}_{\mathfrak{12}}^{\k} = \hat{\bm G}_{0;\mathfrak{12}}^{\k} + \sum_{\mathfrak{ab}}\hat{\bm G}_{0;\mathfrak{1a}}^{\k}\hat{\bm\Sigma}_{\mathfrak{ab}}^{\k}\hat{\bm G}_{\mathfrak{b2}}^{\k},
\end{align}
where
\begin{align}
	\hat{\bm G}_{0;\mathfrak{12}}^{\k}=\begin{pmatrix}
		G_{0;\mathfrak{12}}^{\k} & 0 \\ 0 & G_{0;\mathfrak{12}}^{\k}
	\end{pmatrix}
\end{align}
is the non-interacting generalized propagator and
\begin{align}
	\hat{\bm\Sigma}_{\mathfrak{12}}^{\k}=\begin{pmatrix}
		\Sigma_{\mathfrak{12}}^{\k} & \Delta_{\mathfrak{12}}^{\k} \\ 
		\overline{\Delta}_{\mathfrak{12}}^{\k} & \overline{\Sigma}_{\mathfrak{12}}^{\k}
	\end{pmatrix}
\end{align}
denotes the generalized self-energy, where the diagonal entries $\Sigma_{\mathfrak{12}}^{\k}$ and $\overline{\Sigma}_{\mathfrak{12}}^{\k}$ are the single-particle and single-hole self-energies, respectively and the off-diagonal \textit{pairing} function $\Delta_{\mathfrak{12}}^{\k}$ is coined the superconducting \textit{gap function}. $\Delta_{\mathfrak{12}}^{\k}$ is directly connected to the anomalous propagators and therefore can capture the relevant Cooper-pairing nature of the electrons in the superconducting phase. While the gap function is non-zero only in the superconducting phase, it is nevertheless possible to obtain a solution directly from the normal state: The phase transition from the normal to the superconducting state is characterized by a divergence in the respective particle-particle susceptibility, which can be studied through an eigenvalue problem also known as the Eliashberg equation \cite{Berges2016}. We will start with the BSE written in terms of singlet and triplet susceptibilities in the particle-particle channel, see \eqref{eq:bethe_salpeter_susc_pp},
\begin{align}
\chi^{\q\k\kp}_{\sing/\trip;\mathfrak{1234}}=\chi^{\q\k\kp}_{\sing/\trip;0;\mathfrak{1234}}-\frac{1}{2\beta^2}\sum_{\k_1\k_2;\mathfrak{abcd}}\chi^{\q\k\k_1}_{\sing/\trip;0;\mathfrak{1c3a}}\Gamma^{\q\k_1\k_2}_{\sing/\trip;\mathfrak{abcd}}\left (\chi^{\q\k_2\kp}_{\sing/\trip;\mathfrak{d2b4}}+\chi^{\q\k_2\kp}_{\sing/\trip;0;\mathfrak{d2b4}}\right ),
\end{align}
where the bare bubble susceptibility in the singlet and triplet channel is given by \cite{Gingras2019}
\begin{align}
	\chi^{\q\k\kp}_{\sing/\trip;0;\mathfrak{1234}}=\mp\chi^{\q\k\kp}_{0;\mathfrak{1234}}.
\end{align}
When describing the superconducting state, one usually employs quantities written in the Nambu basis, i.e.\ we are now using the generalized Gro'kov single-particle propagator in the equation for the susceptibility, not the regular Green's function \cite{Nambu1950, Putzer2023}. We will denote these Green's functions by an additional superscript, $G^{\text{sc};\k}_{\mathfrak{12}}$, and call them \textit{superconducting} Green's functions. The corresponding Eliashberg equation can then be written as
\begin{align}\label{eq:full_eliashberg}
	\lambda_{\sing/\trip}\Delta_{\sing/\trip;\mathfrak{12}}^{\k}=\pm\frac{1}{2}\sum_{\k_1\k_2;\mathfrak{abcd}}\Gamma^{\q\k\k_1}_{\sing/\trip;\mathfrak{1b2a}}\chi^{\text{sc};\q\k_1\k_2}_{\sing/\trip;0;\mathfrak{acbd}}\Delta_{\sing/\trip;\mathfrak{dc}}^{\k_2},
\end{align}
yielding an eigenvalue problem with $\lambda_{\sing/\trip}$ as the corresponding eigenvalue. The irreducible vertex $\Gamma^{\q\k\kp}_{\sing/\trip;\mathfrak{1234}}$ appearing in \eqref{eq:full_eliashberg} describes the scattering of two electrons with spin-orbit index $\mathfrak{1},\mathfrak{2}$ and momenta $(\k+\q,-\k)$ to electrons with spin-orbit index $\mathfrak{3},\mathfrak{4}$ and momenta $(\kp+\q,-\kp)$. Here, the bare bubble susceptibility $\chi^{\text{sc};\q\k\kp}_{0;\mathfrak{1234}}$ is now built from the superconducting Green's function $G^{\text{sc};\k}_{\mathfrak{12}}$ and therefore contains the anomalous self-energy in the following form
\begin{align}
	G^{\text{sc};\k}_{\mathfrak{12}}=\frac{i\nu\delta_{\mathfrak{12}}+\varepsilon^{\vv{k}}_{\mathfrak{12}}+\Sigma^{\k}_{\mathfrak{12}}}{(i\nu\delta_{\mathfrak{12}}+\varepsilon^{\vv{k}}_{\mathfrak{12}}+\Sigma^{\k}_{\mathfrak{12}})(i\nu\delta_{\mathfrak{12}}+\varepsilon^{\vv{k}}_{\mathfrak{12}}-\Sigma^{\k}_{\mathfrak{12}})-|\Delta_{\mathfrak{12}}^{\k}|^2},
\end{align}
which results from the generalized Dyson equation of \eqref{eq:grokov_dyson}\footnote{$\Delta_{\sing/\trip;\mathfrak{12}}^{\k}$, as already mentioned above, is the superconducting gap function that transitions smoothly into the anomalous self-energy below the critical temperature $T_\text{c}$.}. A divergence of $\chi^{\q\k\kp}_{\sing/\trip;\mathfrak{1234}}$ occurs when any eigenvalue of $-\Gamma^{\q\k\k_1}_{\sing/\trip;\mathfrak{1b2a}}\chi^{\text{sc};\q\k_1\k_2}_{\sing/\trip;0;\mathfrak{acbd}}$ reaches unity\footnote{This eigenvalue does not necessarily have to be the largest one. Some eigenvalues may exceed unity but remain irrelevant if they are not the ones that first trigger the divergence in the channel under consideration.}. In other words, as the temperature is lowered, the superconducting instability occurs at $T_\text{c}$, where $\lambda_{\sing/\trip}\to 1$\footnote{A common numerical approach to determine $T_\text{c}$ is to compute $\lambda_{\sing/\trip}$ at several low temperatures and extrapolate to the point where $\lambda_{\sing/\trip}\to1$.}. \eqref{eq:full_eliashberg} is the full Eliashberg equation, which represents a non-linear implicit equation for the superconducting gap function $\Delta_{\mathfrak{12}}^{\k}$. To solve this equation numerically, one often linearizes the Eliashberg equation and neglects the gap function in the denominator of the superconducting Green's function. This is then commonly coined the \textit{linearized} Eliashberg equation and reads
\begin{align}\label{eq:linearized_eliashberg}
	\lambda_{\sing/\trip}\Delta_{\sing/\trip;\mathfrak{12}}^{\k}=\pm\frac{1}{2}\sum_{\k_1\k_2;\mathfrak{abcd}}\!\!\Gamma^{\q\k\k_1}_{\sing/\trip;\mathfrak{1b2a}}\chi^{\q\k_1\k_2}_{\sing/\trip;0;\mathfrak{acbd}}\Delta_{\sing/\trip;\mathfrak{dc}}^{\k_2},
\end{align}
where the superconducting Green's function transitions to the regular Green's function for a vanishing $\Delta_{\mathfrak{12}}^{\k}$. \eqref{eq:linearized_eliashberg} is diagrammatically depicted for the case $\q=0$ in \figref{fig:linearized_eliashberg}. Choosing this single point of vanishing transfer momentum simplifies the problem, as Cooper pairs form at zero total energy (i.e.\ in the centre-of-mass system), making the equation effectively uniform and time-independent (c.f. Migdal-Eliashberg theory). Finite $\q$ can indicate competing (pair-)density wave states, i.e.\ charge density waves or spin density waves \cite{Yoshida2021}.
\begin{figure}[ht!]
	\centering
  	\includestandalone[scale=1]{Graphics/Diagrams/linearized_eliashberg/linearized_eliashberg}
  	\caption{Linearized Eliashberg equation for vanishing transfer frequency $\q=0$, see \eqref{eq:linearized_eliashberg}.}
  	\label{fig:linearized_eliashberg}
\end{figure}

Due to Pauli's principle requiring $\Delta$ to be antisymmetric under the exchange of two particles, the combined symmetry operation of (i) spin flip ($\hat S; \sigma\to\overline{\sigma}$), (ii) parity ($\hat P; \vv{k}\to-\vv{k}$), (iii) orbital exchange ($\hat O; n_1 \leftrightarrow n_2$) and (iv) time reversal ($\hat T; \nu\to-\nu$) has to yield \cite{Riseborough2004}
\begin{align}
	\hat S\hat P\hat O\hat T \Delta_{\mathfrak{12}}^{\k} = -\Delta_{\mathfrak{21}}^{-\k}.
\end{align}
The pairing type of the Cooper pairs can be uniquely determined by individually applying the symmetry operations (i)-(iv). The (anti-)symmetric behavior of $\Delta_{\mathfrak{12}}^{\k}$ under these operations of course limits the solutions of the linearized Eliashberg equations to a small subset, which is important to be aware of when solving them numerically. Additionally, the corresponding gap function for singlet pairing fulfills $\Delta^{\k}_{\mathfrak{12}}=\Delta^{-\k}_{\mathfrak{21}}$, whereas the corresponding gap function for triplet pairing fulfills $\Delta^{\k}_{\mathfrak{12}}=-\Delta_{\mathfrak{21}}^{-\k}$ \cite{Nourafkan2016}.

One can solve the linearized Eliashberg equation through a power iteration on $\bm M = -\bm\Gamma\bm\chi$, e.g.\ a Lanczos algorithm. There, one specifies a (normalized) starting gap function $\bm\Delta_0$ and then iterates through the following equations self-consistently:
\begin{subequations}
\begin{align}
	\bm\Delta_{n} &= \bm M\bm\Delta_{n-1},\\
	\bm\Delta_{n} &= \frac{\bm\Delta_{n}}{||\bm\Delta_{n}||}\qqand\\
	\lambda_{n} &= \bm\Delta_{n}^{\dagger}\bm M\bm\Delta_{n},	
\end{align}
\end{subequations}
where $n$ denotes the iteration number. Solving this set of equations is not difficult, the computationally demanding part of solving the linearized Eliashberg equation is the calculation of the irreducible vertex in the particle-particle channel. Following the derivation outlined in the supplemental material of \cite{Kitatani2019} and generalizing the equations to account for multiple orbitals yields for the full vertex in ladder-$\dga$ for $\r\in\{\dens/\magn\}$ spin combinations
\begin{align}\label{eq:full_vertex_ladder_dga_non_crossing_symmetric}
\begin{split}
	F^{\q\nu\nu'}_{\r;\mathfrak{1234}}&=\beta^2\Big [\left (\chi_{0;\mathfrak{1234}}^{\q\nu\nu'}\right )^{-1}-\!\!\sum_{\nu_1\nu_2;\mathfrak{abcd}}\!\!\left (\chi_{0;\mathfrak{12ba}}^{\q\nu\nu_1}\right )^{-1} \chi_{\r;\mathfrak{abcd}}^{*;\q\nu_1\nu_2} \left (\chi_{0;\mathfrak{dc34}}^{\q\nu_2\nu'}\right )^{-1}\Big ]\\
	&\qquad\qquad + \sum_{\substack{\mathfrak{abcd}\\\mathfrak{efgh}}} \gamma^{\q\nu}_{\r;\mathfrak{12ab}} \mathcal{U}^{\vv{q}}_{r;\mathfrak{badc}}\left (\unity_{\mathfrak{cdgh}}-\mathcal{U}^{\vv{q}}_{\r;\mathfrak{cdef}}\chi^{\q}_{\r;\mathfrak{fegh}}\right )\gamma^{\q\nu'}_{\r;\mathfrak{hg34}}, 
\end{split}
\end{align}
with the auxiliary susceptibility $\chi_{\r;\mathfrak{1234}}^{*;\q\nu\nu'}$ of \eqref{eq:chi_aux} and the three-leg vertex $\gamma^{\q\nu}_{\r;\mathfrak{1234}}$ of \eqref{eq:vrg}. In the following section, we will present a short derivation of the particle-particle pairing vertex for the case of $\q=0$ within the framework of ladder-$\dga$. 

\subsection{Derivation of the pairing vertex $\Gamma_{\pp}$ within ladder-$\dga$}

The vertex in \eqref{eq:full_vertex_ladder_dga_non_crossing_symmetric} does not fulfill any crossing-symmetry relation, as it is generated solely from the particle-hole BSE\footnote{It is solely the ladder-vertex and not the full vertex in the ladder-$\dga$ formalism.}. Consequently, it only possesses particle-swapping symmetry \cite{Galler2017}. In contrast, the full vertex in ladder-$\dga$, given in \eqref{eq:full_vertex_ladder_dga_crossing_symmetric}, satisfies crossing symmetry. For the $\uparrow\downarrow$-component, we have \cite{Galler2017, Rohringer2013}
\begin{align}
	F^{\dga;\q\k\kp}_{\mathfrak{1234};\uparrow\downarrow}&=\frac12\left (F_{\dens;\mathfrak{1234}}^{\q\nu\nu'}-F_{\magn;\mathfrak{1234}}^{\q\nu\nu'}\right )-F_{\magn;\mathfrak{1432}}^{(\k-\kp)\nu(\nu-\omega)}-\frac12\left (F_{\dens;\mathfrak{1234}}^{\omega\nu\nu'}-F_{\magn;\mathfrak{1234}}^{\omega\nu\nu'}\right ).
\end{align}
The $F_{\magn;\mathfrak{1432}}^{(\k-\kp)\nu(\nu-\omega)}$-term arises from the $\ph\to\phbar$ transformation that restores crossing symmetry, introducing an orbital permutation and frequency shifts. This expression is written only for the $\uparrow\downarrow$-component. Since we are interested in the singlet ($\uparrow\downarrow-\,\overline{\uparrow\downarrow}$) and triplet ($\uparrow\downarrow+\,\overline{\uparrow\downarrow}$) spin combinations, the $\overline{\uparrow\downarrow}$-component is also needed. Fortunately, one can obtain the $\overline{\uparrow\downarrow}$-contribution via a crossing symmetry relation of the $\uparrow\downarrow$-component in the following way:
\begin{align}\label{eq:f_ud_bar_to_ud}
	F^{\q\k\kp}_{\mathfrak{1234};\overline{\uparrow\downarrow}}=-F^{(\k-\kp)\k(\k-\q)}_{\mathfrak{1432};\uparrow\downarrow}.
\end{align}
One therefore obtains for the $\overline{\uparrow\downarrow}$-component via \eqref{eq:f_ud_bar_to_ud}
\begin{align}
\begin{split}
	F^{\dga;\q\k\kp}_{\mathfrak{1234};\overline{\uparrow\downarrow}}&=-\frac12\left (F_{\dens;\mathfrak{1432}}^{(\k-\kp)\nu(\nu-\omega)}-F_{\magn;\mathfrak{1432}}^{(\k-\kp)\nu(\nu-\omega)}\right )+F_{\magn;\mathfrak{1234}}^{\q\nu\nu'}\\
	&\qquad+\frac12\left (F_{\dens;\mathfrak{1432}}^{(\nu-\nu')\nu(\nu-\omega)}-F_{\magn;\mathfrak{1432}}^{(\nu-\nu')\nu(\nu-\omega)}\right ).
\end{split}
\end{align}
This yields for the full vertices in $\ph$-notation for both the singlet
\begin{align}
\begin{split}
	F_{\sing;\mathfrak{1234}}^{\dga;\q\k\kp}&=\frac12\left (F_{\dens;\mathfrak{1234}}^{\q\nu\nu'}-3F_{\magn;\mathfrak{1234}}^{\q\nu\nu'}\right )-\frac12\left (F_{\dens;\mathfrak{1234}}^{\omega\nu\nu'}-F_{\magn;\mathfrak{1234}}^{\omega\nu\nu'}\right )\\
	&\qquad+\frac12\left (F_{\dens;\mathfrak{1432}}^{(\k-\kp)\nu(\nu-\omega)}-3F_{\magn;\mathfrak{1432}}^{(\k-\kp)\nu(\nu-\omega)}\right )-\frac12\left (F_{\dens;\mathfrak{1432}}^{(\nu-\nu')\nu(\nu-\omega)}-F_{\magn;\mathfrak{1432}}^{(\nu-\nu')\nu(\nu-\omega)}\right )
\end{split}
\end{align}
and triplet spin combination
\begin{align}
\begin{split}
	F_{\trip;\mathfrak{1234}}^{\dga;\q\k\kp}&=\frac12\left (F_{\dens;\mathfrak{1234}}^{\q\nu\nu'}+F_{\magn;\mathfrak{1234}}^{\q\nu\nu'}\right )-\frac12\left (F_{\dens;\mathfrak{1234}}^{\omega\nu\nu'}-F_{\magn;\mathfrak{1234}}^{\omega\nu\nu'}\right )\\
	&\qquad-\frac12\left (F_{\dens;\mathfrak{1432}}^{(\k-\kp)\nu(\nu-\omega)}+F_{\magn;\mathfrak{1432}}^{(\k-\kp)\nu(\nu-\omega)}\right )+\frac12\left (F_{\dens;\mathfrak{1432}}^{(\nu-\nu')\nu(\nu-\omega)}-F_{\magn;\mathfrak{1432}}^{(\nu-\nu')\nu(\nu-\omega)}\right ).
\end{split}
\end{align}
We now want to evaluate this expression for $\q_{\pp}=0$ by setting $\q_{\ph}=\k_{\pp}+\kp_{\pp}, \k_{\ph}=\k_{\pp}$ and $\kp_{\ph}=\kp_{\pp}$\footnote{The frequency notation for the particle-particle channel used in this thesis has been written down in \eqref{eq:pp_freq}.}. Inserting these frequencies and momenta yields for the singlet channel
\begin{align}\label{eq:ladder_f_sing}
\begin{split}
	F_{\sing;\mathfrak{1234}}^{\dga;(\q_{\pp}=0)\k\kp}&=\frac12\left (F_{\dens;\mathfrak{1234}}^{(\k+\kp)\nu\nu'}-3F_{\magn;\mathfrak{1234}}^{(\k+\kp)\nu\nu'}\right )-\frac12\left (F_{\dens;\mathfrak{1234}}^{(\nu+\nu')\nu\nu'}-F_{\magn;\mathfrak{1234}}^{(\nu+\nu')\nu\nu'}\right )\\
	&\qquad+\frac12\left (F_{\dens;\mathfrak{1432}}^{(\k-\kp)\nu(-\nu')}-3F_{\magn;\mathfrak{1432}}^{(\k-\kp)\nu(-\nu')}\right )-\frac12\left (F_{\dens;\mathfrak{1432}}^{(\nu-\nu')\nu(-\nu')}-F_{\magn;\mathfrak{1432}}^{(\nu-\nu')\nu(-\nu')}\right )\\
	&=\underbrace{\frac12\left (F_{\dens;\mathfrak{1234}}^{(\k+\kp)\nu\nu'}-3F_{\magn;\mathfrak{1234}}^{(\k+\kp)\nu\nu'}\right )-F_{\mathfrak{1234};\uparrow\downarrow}^{(\nu+\nu')\nu\nu'}}_{\textstyle\tilde{F}_{\sing;\mathfrak{1234}}^{\dga;(\q=0)\k\kp}}+ \underbrace{\left (\kp\rightarrow -\kp\;\&\; \mathfrak{1234}\rightarrow\mathfrak{1432}\right )}_{\textstyle\tilde{F}_{\sing;\mathfrak{1432}}^{\dga;(\q=0)\k(-\kp)}}.
\end{split}
\end{align}
For the triplet channel we analogously find
\begin{align}\label{eq:ladder_f_trip}
\begin{split}
	F_{\trip;\mathfrak{1234}}^{\dga;(\q_{\pp}=0)\k\kp}&=\frac12\left (F_{\dens;\mathfrak{1234}}^{(\k+\kp)\nu\nu'}+F_{\magn;\mathfrak{1234}}^{(\k+\kp)\nu\nu'}\right )-\frac12\left (F_{\dens;\mathfrak{1234}}^{(\nu+\nu')\nu\nu'}-F_{\magn;\mathfrak{1234}}^{(\nu+\nu')\nu\nu'}\right )\\
	&\qquad-\frac12\left (F_{\dens;\mathfrak{1432}}^{(\k-\kp)\nu(-\nu')}+F_{\magn;\mathfrak{1432}}^{(\k-\kp)\nu(-\nu')}\right )+\frac12\left (F_{\dens;\mathfrak{1432}}^{(\nu-\nu')\nu(-\nu')}-F_{\magn;\mathfrak{1432}}^{(\nu-\nu')\nu(-\nu')}\right )\\
	&=\underbrace{\frac12\left (F_{\dens;\mathfrak{1234}}^{(\k+\kp)\nu\nu'}+F_{\magn;\mathfrak{1234}}^{(\k+\kp)\nu\nu'}\right )-F_{\mathfrak{1234};\uparrow\downarrow}^{(\nu+\nu')\nu\nu'}}_{\textstyle\tilde{F}_{\trip;\mathfrak{1234}}^{\dga;(\q=0)\k\kp}}-\underbrace{\left (\kp\rightarrow -\kp\;\&\; \mathfrak{1234}\rightarrow\mathfrak{1432}\right )}_{\textstyle\tilde{F}_{\trip;\mathfrak{1432}}^{\dga;(\q=0)\k(-\kp)}}.
\end{split}
\end{align}
We can construct the pairing vertex in either singlet or triplet channel by subtracting the purely \textit{local} particle-particle reducible diagrams. We only need to consider the local diagrams here, since neglecting the momentum dependence of particle-particle reducible diagrams is a key approximation in ladder-$\dga$. This yields
\begin{align}
	\Gamma_{\sing/\trip;\mathfrak{1234}}^{(\q=0)\k\kp}=F_{\sing/\trip;\mathfrak{1234}}^{\dga;(\q=0)\k\kp}-\phi_{\sing/\trip;\mathfrak{1234}}^{(\omega=0)\nu\nu'}.
\end{align}
The local reducible diagrams can be subtracted straightforwardly from $\tilde{F}_{\sing/\trip;\mathfrak{1234}}^{\dga;(\q_{\pp}=0)\k(\pm\kp)}$ in the following fashion
\begin{align}\label{eq:pairing_vertex_full_formula}
\begin{split}
	\Gamma_{\sing/\trip;\mathfrak{1234}}^{(\q=0)\k\kp}&=\tilde{F}_{\sing/\trip;\mathfrak{1234}}^{\dga;(\q=0)\k\kp}\pm \tilde{F}_{\sing/\trip;\mathfrak{1432}}^{\dga;(\q=0)\k(-\kp)}-\phi_{\sing/\trip;\mathfrak{1234}}^{(\omega=0)\nu\nu'}\\
	&=\tilde{F}_{\sing/\trip;\mathfrak{1234}}^{\dga;(\q=0)\k\kp}\pm \tilde{F}_{\sing/\trip;\mathfrak{1432}}^{\dga;(\q=0)\k(-\kp)}-\left (\phi_{\pp;\mathfrak{1234};\uparrow\downarrow}^{(\omega=0)\nu\nu'}\mp \phi_{\pp;\mathfrak{1234};\overline{\uparrow\downarrow}}^{(\omega=0)\nu\nu'}\right )\\
	&=\tilde{F}_{\sing/\trip;\mathfrak{1234}}^{\dga;(\q=0)\k\kp}\pm \tilde{F}_{\sing/\trip;\mathfrak{1432}}^{\dga;(\q=0)\k(-\kp)}-\left (\phi_{\pp;\mathfrak{1234};\uparrow\downarrow}^{(\omega=0)\nu\nu'}\pm \phi_{\pp;\mathfrak{1432};\uparrow\downarrow}^{(\omega=0)\nu(-\nu')}\right )\\
	&=\underbrace{\tilde{F}_{\sing/\trip;\mathfrak{1234}}^{\dga;(\q=0)\k\kp}-\phi_{\pp;\mathfrak{1234};\uparrow\downarrow}^{(\omega=0)\nu\nu'}}_{\textstyle\tilde{\Gamma}_{\sing/\trip;\mathfrak{1234}}^{(\q=0)\k\kp}}\pm \underbrace{\left (\tilde{F}_{\sing/\trip;\mathfrak{1432}}^{\dga;(\q=0)\k(-\kp)}- \phi_{\pp;\mathfrak{1432};\uparrow\downarrow}^{(\omega=0)\nu(-\nu')}\right )}_{\textstyle\tilde{\Gamma}_{\sing/\trip;\mathfrak{1432}}^{(\q=0)\k(-\kp)}}\\
	&=\tilde{\Gamma}_{\sing/\trip;\mathfrak{1234}}^{(\q=0)\k\kp}\pm \left (\kp\rightarrow -\kp\;\&\; \mathfrak{1234}\rightarrow\mathfrak{1432}\right ).
\end{split}
\end{align}

\subsection{Symmetry of the gap function}

We can show that the gap function inherits the correct symmetry upon using the symmetrized vertices of \eqref{eq:pairing_vertex_full_formula}. Recall that the multi-orbital Eliashberg equation for the case $\q=0$ is given by 
\begin{align}
	\lambda_{\sing/\trip}\Delta_{\sing/\trip;\mathfrak{12}}^{\k}=\pm\frac{1}{2}\sum_{\kp;\mathfrak{abcd}}\!\!\Gamma^{(\q=0)\k\kp}_{\sing/\trip;\mathfrak{1b2a}}\chi^{(\q=0)\kp}_{\sing/\trip;0;\mathfrak{acbd}}\Delta_{\sing/\trip;\mathfrak{dc}}^{\kp}
\end{align}
for $\q=0$ and the singlet and triplet gap functions should fulfill $\Delta_{\sing;\mathfrak{12}}^{\k}=\Delta_{\sing;\mathfrak{21}}^{-\k}$ and $\Delta_{\trip;\mathfrak{12}}^{\k}=-\Delta_{\trip;\mathfrak{21}}^{-\k}$, respectively. We can write $\Gamma$ in the symmetrized form with the irreducible vertices $\tilde{\Gamma}$ from \eqref{eq:pairing_vertex_full_formula} above, i.e.\
\begin{align}\label{eq:symmetrization_pairing_vertex}
	\Gamma^{(\q=0)\k\kp}_{\sing/\trip;\mathfrak{1234}} = \tilde{\Gamma}^{(\q=0)\k\kp}_{\sing/\trip;\mathfrak{1234}}\pm \tilde{\Gamma}^{(\q=0)\k(-\kp)}_{\sing/\trip;\mathfrak{1432}},
\end{align}
where $\tilde{\Gamma}^{(\q=0)\k\kp}_{\sing/\trip;\mathfrak{1234}}$ contains both the non-local and local contribution to the irreducible vertex. To confirm that this combination enforces the correct symmetry of the gap function, we need to show that
\begin{subequations}
\begin{align}
	\Delta_{\sing;\mathfrak{12}}^{\k}-\Delta_{\sing;\mathfrak{21}}^{-\k}&=0\qqand\\
	\Delta_{\trip;\mathfrak{12}}^{\k}+\Delta_{\trip;\mathfrak{21}}^{-\k}&=0.
\end{align}
\end{subequations}
This expression for $\q=0$ then reads for both the singlet and triplet channels
\begin{subequations}
\begin{align}
	\Delta_{\sing;\mathfrak{12}}^{\k}-\Delta_{\sing;\mathfrak{21}}^{-\k}&\propto\!\!\sum_{\kp;\mathfrak{abcd}}\!\! \left [\Gamma^{(\q=0)\k\kp}_{\sing;\mathfrak{1b2a}}-\Gamma^{(\q=0)(-\k)\kp}_{\sing;\mathfrak{2b1a}}\right ]\chi^{(\q=0)\kp}_{\sing;0;\mathfrak{acbd}}\Delta_{\sing;\mathfrak{dc}}^{\kp}\label{eq:symmetrization_delta_first_formula_singlet}\qqand\\
	\Delta_{\trip;\mathfrak{12}}^{\k}+\Delta_{\trip;\mathfrak{21}}^{-\k}&\propto\!\!\sum_{\kp;\mathfrak{abcd}}\!\! \left [\Gamma^{(\q=0)\k\kp}_{\trip;\mathfrak{1b2a}}+\Gamma^{(\q=0)(-\k)\kp}_{\trip;\mathfrak{2b1a}}\right ]\chi^{(\q=0)\kp}_{\trip;0;\mathfrak{acbd}}\Delta_{\trip;\mathfrak{dc}}^{\kp}\label{eq:symmetrization_delta_first_formula_triplet}.
\end{align}
\end{subequations}
The full non-local ladder vertices of \eqref{eq:ladder_f_sing} and \eqref{eq:ladder_f_trip} fulfill the swapping symmetry, which is passed along to the pairing vertex $\Gamma^{(\q=0)\k\kp}_{\sing/\trip;\mathfrak{1234}}$ and reads in $\pp$-notation for $\q=0$
\begin{align}\label{eq:swapping_symmetry_gamma}
	\Gamma^{(\q=0)\k\kp}_{\sing/\trip;\mathfrak{1234}}=\Gamma^{(\q=0)(-\k)(-\kp)}_{\sing/\trip;\mathfrak{3412}}\overset{\text{(\ref{eq:symmetrization_pairing_vertex})}}{=}\pm \Gamma^{(\q=0)(-\k)\kp}_{\sing/\trip;\mathfrak{3214}}.
\end{align}
Applying \eqref{eq:swapping_symmetry_gamma} to \eqref{eq:symmetrization_delta_first_formula_singlet} then yields for the singlet channel
\begin{align}
	\Delta_{\sing;\mathfrak{12}}^{\k}-\Delta_{\sing;\mathfrak{21}}^{-\k}&\propto\!\!\sum_{\kp;\mathfrak{abcd}}\!\! \left [\Gamma^{(\q=0)\k\kp}_{\sing;\mathfrak{1b2a}}-\Gamma^{(\q=0)\k\kp}_{\sing;\mathfrak{1b2a}}\right ]\chi^{(\q=0)\kp}_{\sing;0;\mathfrak{acbd}}\Delta_{\sing;\mathfrak{dc}}^{\kp}=0
\end{align}
and analogously for the triplet channel, see \eqref{eq:symmetrization_delta_first_formula_triplet},
\begin{align}
	\Delta_{\trip;\mathfrak{12}}^{\k}+\Delta_{\trip;\mathfrak{21}}^{-\k}&\propto\!\!\sum_{\kp;\mathfrak{abcd}}\!\! \left [\Gamma^{(\q=0)\k\kp}_{\trip;\mathfrak{1b2a}}-\Gamma^{(\q=0)\k\kp}_{\trip;\mathfrak{1b2a}}\right ]\chi^{(\q=0)\kp}_{\trip;0;\mathfrak{acbd}}\Delta_{\trip;\mathfrak{dc}}^{\kp}=0,
\end{align}
confirming that the symmetrization in \eqref{eq:symmetrization_pairing_vertex} leads to the correct symmetry of the gap function.

\subsection{Efficient solving of the Eliashberg equation via a Fourier transform}

To avoid having to deal with two momentum indices for the pairing vertex $\Gamma_{\mathfrak{1234}}^{(\q=0)\k\k'}$ in the calculation, one can perform a Fourier transform to real space. However, this is only of advantage in the ladder approximation of the vertex, since then $\Gamma_{\sing/\trip;\mathfrak{1234}}^{(\q=0)\k\kp}=\Gamma_{\sing/\trip;\mathfrak{1234}}^{\vv{k}-\vv{k}',(\omega=0)\nu\nu'}\pm\Gamma_{\sing/\trip;\mathfrak{1432}}^{\vv{k}+\vv{k}',(\omega=0)\nu(-\nu')}$. The sum over $\k'$ of \eqref{eq:linearized_eliashberg} is essentially a convolution in momentum space, which reduces to a multiplication in Fourier space. We thus take the Fourier transform of \eqref{eq:linearized_eliashberg} in the case of $\q=0$ with respect to $\vv{k}$
\begin{align}\label{eq:eliashberg_fourier_transform_partial}
	\lambda_{\sing/\trip}\sum_{\vv{k}}\text{e}^{\text{i}\vv{kx}}\;\Delta_{\sing/\trip;\mathfrak{12}}^{\k}=\pm\frac{1}{2}\sum_{\vv{k}\kp;\mathfrak{abcd}}\!\!\text{e}^{\text{i}\vv{kx}}\;\left [\Gamma_{\sing/\trip;\mathfrak{1b2a}}^{\vv{k}-\vv{k}'(\omega=0)\nu\nu'}\pm\Gamma_{\sing/\trip;\mathfrak{1a2b}}^{\vv{k}+\vv{k}'(\omega=0)\nu(-\nu')}\right ]\chi^{\kp}_{\sing/\trip;0;\mathfrak{acbd}}\Delta_{\sing/\trip;\mathfrak{dc}}^{\kp},
\end{align}
where $\vv{x}=(x_1,...,x_D)$ is a real-space vector for systems in $D$ dimensions. If we introduce
\begin{align}
	\tilde\Delta^{\kp}_{\sing/\trip;\mathfrak{ba}}=\sum_{\mathfrak{cd}}\chi^{\kp}_{\sing/\trip;0;\mathfrak{acbd}}\Delta_{\sing/\trip;\mathfrak{dc}}^{\kp}
\end{align}
and use the shifts $\vv{k}=\vv{q}+\vv{k}'$ for the left term and $\vv{k}=\vv{q}-\vv{k}'$ for the right term, we can rewrite \eqref{eq:eliashberg_fourier_transform_partial} as
\begin{align}
	\lambda_{\sing/\trip}\sum_{\vv{k}}\text{e}^{\text{i}\vv{kx}}\; \Delta_{\sing/\trip;\mathfrak{12}}^{\k}=\pm\frac12\sum_{\nu'\mathfrak{ab}}\sum_{\vv{q}\vv{k}'}\text{e}^{\text{i}\vv{qx}}\;\left [\Gamma_{\sing/\trip;\mathfrak{1b2a}}^{\vv{q}(\omega=0)\nu\nu'}\text{e}^{\text{i}\vv{k}'\vv{x}}\;\tilde\Delta^{\kp}_{\sing/\trip;\mathfrak{ba}}\pm\Gamma_{\sing/\trip;\mathfrak{1a2b}}^{\vv{q}(\omega=0)\nu(-\nu')}\text{e}^{-\text{i}\vv{k}'\vv{x}}\;\tilde\Delta^{-\vv{k}'\nu'}_{\sing/\trip;\mathfrak{ba}}\right ].
\end{align}
This reads in Fourier-transformed quantities as
\begin{align}
	\lambda_{\sing/\trip} \Delta^{\vv{x}\nu}=\pm\frac12\sum_{\nu'\mathfrak{ab}}\left [\Gamma^{\vv{x}(\omega=0)\nu\nu'}_{\sing/\trip;\mathfrak{1b2a}}\tilde\Delta^{\vv{x}\nu'}_{\sing/\trip;\mathfrak{ba}}\pm\Gamma^{-\vv{x}(\omega=0)\nu(-\nu')}_{\sing/\trip;\mathfrak{1a2b}}\tilde\Delta^{-\vv{x}\nu'}_{\sing/\trip;\mathfrak{ba}}\right ].
\end{align}
In essence, we have reduced the convolution of \eqref{eq:linearized_eliashberg} in momentum space to a product of the quantities in real space, effectively reducing the object from two momentum dimensions $N_{\vv{k}}\times N_{\vv{k}}$ to a single real-space dimension $N_{\vv{x}}(=N_{\vv{k}})$, reducing the object size by a factor $N_{\vv{k}}$. Numerically, solving the equation now can be done by performing a discrete Fourier transform on the quantities before executing the power iteration to calculate the gap function on the real-space axis. To afterwards obtain the momentum-dependent gap function, one simply applies the inverse transformation to the result.

\end{document}
